\begin{exercisechunk}
\item \textbf{Symbolic histories of the tent map.}
\label[exercise]{ex:l14-tent-symbolic-dynamics}
\exerciseprofile{2}{2}{2}{\exproof\quad\excalculation\quad\exinterpretation}{%
Connect the affine pieces of an iterated tent map with uniformly distributed
binary symbolic histories.}
The term \emph{symbolic dynamics} records which branch of a dynamical system
an orbit visits. This exercise shows that the \(2^L\) affine pieces of the
tent-map iterate correspond exactly to \(L\)-bit branch histories.

Let \(X\) be uniformly distributed on \([0,1]\) under the ambient probability
measure \(\P\). Define
\[
 B_t=\I\bigl(\tau^{\circ t}(X)\ge1/2\bigr),
 \qquad t=0,1,\ldots.
\]
Events involving dyadic endpoints may be assigned to either adjacent branch;
these endpoints have probability zero.
\begin{enumerate}[(a),beginpenalty=10000]
\item Prove directly from the two inverse branches of \(\tau\) that
\[
 \P(\tau(X)\le y)=y,
 \qquad 0\le y\le1.
\]
Thus \(\tau(X)\) is uniformly distributed on \([0,1]\).
\item For \(b=(b_0,\ldots,b_{L-1})\in\{0,1\}^L\), define
\[
 C_b=\left\{x\in[0,1]:
 \I(\tau^{\circ t}(x)\ge1/2)=b_t
 \text{ for }t=0,\ldots,L-1\right\}.
\]
Prove by induction on \(L\) that, up to endpoints, \(C_b\) is one interval
of length \(2^{-L}\), and \(\tau^{\circ L}\) maps that interval affinely
onto \([0,1]\).
\item Deduce that
\[
 \P(B_0=b_0,\ldots,B_{L-1}=b_{L-1})=2^{-L}
 \qquad(b\in\{0,1\}^L).
\]
Conclude that \(B_0,\ldots,B_{L-1}\) are independent Bernoulli random
variables with success probability \(1/2\).
\item Use part~(a) inductively to identify the law of
\(\tau^{\circ t}(X)\) for every \(t\ge0\). Explain how parts~(b)--(c)
relate the \(2^L\) affine pieces to symbolic histories rather than merely
counting pieces.
\end{enumerate}
\end{exercisechunk}
